\subsection{Compiling a reference configuration}
\label{sec:tutor-reference}

Compilation of \NEMOSIcu{} is done with the script \filename{makenemo} (no extension) in the top-level source directory.
Platform-specific compiler options are specified in a configuration file in the \dirname{arch} subdirectory.
This is a similar workflow to XIOS (\refsection{sec:tutor-xios}), but for NEMO there is no \filename{.path} file.
The \filename{.env} and \filename{.fcm} files for compiling on RACC2 can be found in the auxiliary resources under \dirname{nemo/arch}.

\begin{enumerate}
    \item Make a copy of the two files under the \dirname{arch} subdirectory of your top-level NEMO source directory.
\end{enumerate}

For best results, \NEMOSIcu{} should be compiled to use the XIOS library for input and output as described in \refsection{sec:tutor-xios}.
Assuming that has been done, your copy of the configuration file for NEMO must be edited so that NEMO knows where to find your compiled XIOS library:

\begin{enumerate}\setcounter{enumi}{1}
    \item Open your \filename{arch-RACC2.env} file and, if required, change the value of the environment variable \texttt{XIOS_HOME} to the path to your XIOS installation.

    The placeholder for this, \texttt{\$\{HOME\}/software/XIOS2}, corresponds to the location of XIOS2 used in \refsection{sec:tutor-xios}, so if you used the same path then this step can be skipped.
\end{enumerate}

To compile \NEMOSIcu{} requires selecting a \textit{configuration} of model components (and other pre-compilation options).
A set of reference configurations is provided under the \dirname{cfgs} subdirectory of the top-level source directory.
To begin, we shall compile the \texttt{ORCA2_ICE_PISCES} reference configuration.
This includes the full NEMO ocean model on the ORCA2 (nominal resolution 2\textdegree{}) global grid, the \SIcu{} sea ice model, and the PISCES biogeochemical model.
We shall not run this configuration because PISCES turns out to be a resource hog and is not needed for most sea ice work.
However, compiling this standard configuration will act as a check that the build machinery is working properly, as well as introduce the basic compilation command, before we get on to custom configurations in \refsection{sec:tutor-custom}.

\begin{enumerate}\setcounter{enumi}{2}
    \item Compile the \texttt{ORCA2_ICE_PISCES} configuration with the following command:
\begin{VerbatimCMD}
./makenemo \cmdflag{-m} RACC2 \cmdflag{-r} ORCA2_ICE_PISCES
\end{VerbatimCMD}
\end{enumerate}

Compilation will take a few minutes, during which time \filename{makenemo} gives a commentary on what it is doing.
If all goes well, the final lines should be something like:

\begingroup\small
\begin{Verbatim}[samepage=true]
------------------------------------------------------------------------
Compilation successful
------------------------------------------------------------------------

"./nemo" -> "../BLD/bin/nemo.exe"

------------------------------------------------------------------------
Created symbolic link ./nemo to NEMO executable in experiment directory
/home/users/ab123456/nemo/cfgs/ORCA2_ICE_PISCES/EXP00/nemo
------------------------------------------------------------------------
\end{Verbatim}
\endgroup